\chapter{The Cryptography Behind Blockchains}

A blockchain is not a cryptographic primitive. It is a database --- an append-only ledger --- that a crowd of mutually suspicious strangers maintains without a referee. What lets that crowd agree on a single history, and lets any reader check that history without trusting anyone, is cryptography: hash functions chain the records together, Merkle trees let one number stand in for millions of transactions, digital signatures prove who authorized what, and a hash puzzle (proof of work) makes rewriting the past expensive. Almost every tool in this book shows up. Nothing here is new; what is new is the \emph{assembly}.

This chapter takes the primitives apart and shows how they snap together. We will hash a tiny block by hand, build a four-leaf Merkle tree and verify one leaf with a short proof, derive a Bitcoin-style address from a public key, sign a toy transaction with ECDSA, and mine a block against a deliberately easy target so you can watch the nonce search happen. Along the way we will be exact about a point that gets lost in hype: the cryptography secures \emph{integrity and authorization}, and nothing else. It does not make a coin valuable, a contract legal, or a key un-loseable. Confusing those is how people lose money.

\chapterline

\section{The job: agreement without a referee}
\label{ch36:job}

Start with the Alice/Bob/Eve story, adapted. Alice wants to pay Bob ten coins. In a bank, a trusted server debits Alice's balance and credits Bob's; the server's word is final. Take the server away and replace it with a public network anyone can join, including Eve, who would happily spend the same coin twice or rewrite yesterday's payments. We need a way for everyone to agree on \emph{one} ordered list of transactions --- the ledger --- such that:

\begin{itemize}
  \item only the owner of a coin can spend it (\emph{authorization}),
  \item once a transaction is buried in the history, changing it is infeasible (\emph{integrity / immutability}),
  \item any newcomer can download the ledger and verify it from scratch, trusting no single participant (\emph{public verifiability}).
\end{itemize}

\begin{intuitionbox}
Picture a shared notebook passed around a table of people who do not trust each other. Three rules make it work. (1) Each entry is signed by its author with a signature nobody can forge, so you cannot spend my coins. (2) Each page carries a fingerprint of the previous page, so tearing out or editing an old page changes every fingerprint after it --- the tampering is loud, not silent. (3) Writing a new page costs real effort (solving a puzzle), so an attacker who wants to rewrite history must redo all that effort faster than the rest of the table combined. None of these three is a new invention; they are signatures, hash chaining, and proof of work. The blockchain is the way they interlock.
\end{intuitionbox}

The cryptography delivers (1) and (2) outright. Point (3), and the ``faster than everyone combined'' part, is \emph{consensus} --- an economic and game-theoretic argument layered on top of the crypto, not a theorem about hash functions. Keep that seam in mind; we will return to it when we separate what the math guarantees from what it does not.

\begin{historybox}
The chaining-by-hash idea predates Bitcoin by nearly two decades. In 1990--91 Stuart Haber and W.~Scott Stornetta published a scheme for timestamping digital documents by linking each document's hash to the previous one, so that the position of a record in the chain fixed its time of creation; they later folded in Merkle trees (Ralph Merkle's hash-tree patent dates to 1979) to batch many documents per link. In 2008 a pseudonymous author writing as Satoshi Nakamoto combined hash chaining, Merkle trees, ECDSA signatures, and a proof-of-work consensus rule (echoing Adam Back's 1997 Hashcash) into Bitcoin, whose whitepaper cites Haber and Stornetta directly. The novelty was not any single primitive but the consensus rule that let an open, permissionless set of participants agree on which chain is the real one. That design sparked long-running debates over proof-of-work's energy cost versus proof-of-stake, and over how to scale throughput without making full-node verification too expensive for ordinary users.
\end{historybox}

\section{Hashes: the glue}
\label{ch36:hashes}

Everything rests on a collision-resistant hash function (Chapter 14). Bitcoin uses SHA-256, almost always doubled --- $\Hash(x) = \mathrm{SHA256}(\mathrm{SHA256}(x))$, written SHA-256d --- and RIPEMD-160 for addresses; Ethereum uses Keccak-256. The exact choice matters less than the properties we lean on. Recall them, because each blockchain mechanism relies on a specific one:

\begin{itemize}
  \item \textbf{Preimage resistance}: given a digest $y$, you cannot find any $x$ with $\Hash(x)=y$. This is what makes proof of work a search rather than an inversion.
  \item \textbf{Second-preimage resistance}: given $x$, you cannot find $x'\ne x$ with the same digest. This is what stops you from swapping a block's contents while keeping its fingerprint.
  \item \textbf{Collision resistance}: you cannot find \emph{any} pair $x\ne x'$ that collide. Merkle trees and commitments need this; a collision would let someone open a commitment two ways.
\end{itemize}

A hash is also, informally, a \emph{binding commitment}: publishing $\Hash(x)$ pins down $x$ without revealing it, and you cannot later claim it was some other $x'$. That single sentence is the engine of the whole construction. A block header commits to its transactions; a chain of headers commits to all of history; an address commits to a public key.

\begin{connectionbox}
Two facts from Chapter~14 carry the load here. First, the birthday bound: finding a collision in an $n$-bit hash costs about $2^{n/2}$ work, not $2^{n}$, which is why blockchains use 256-bit digests --- $2^{128}$ is the security floor, and that is why a 160-bit address is considered marginal but a 128-bit one would be reckless. Second, modular-exponentiation and elliptic-curve costs from the algorithms companion: signing and verification are a handful of curve operations, cheap enough to do per transaction, while a brute-force key search is $2^{128}$ curve operations, infeasible. The whole edifice is a bet that $2^{128}$ is and stays out of reach.
\end{connectionbox}

\section{Binding a block to its predecessor}
\label{ch36:chain}

A \emph{block} bundles a batch of transactions plus a small \emph{header}. The header is the part that gets hashed and chained. A simplified Bitcoin header contains:

\begin{itemize}
  \item the hash of the \emph{previous} block's header (\texttt{prev}),
  \item the \emph{Merkle root} of this block's transactions (\texttt{root}, Section~\ref{ch36:merkle}),
  \item a \emph{timestamp},
  \item the difficulty \emph{target} (Section~\ref{ch36:pow}),
  \item a \emph{nonce} --- the number miners vary to solve the puzzle.
\end{itemize}

The block's identity is $\mathrm{ID} = \Hash(\text{header})$. Because the header includes \texttt{prev}, each block's ID depends on the one before it, and that one on the one before \emph{it}, all the way back to the first block (the \emph{genesis} block, which has no predecessor and hardcodes \texttt{prev} to all zeros).

\begin{figure}[H]
\centering
\begin{tikzpicture}[font=\small,
  blk/.style={draw=sourceblue!80, rounded corners=2pt, fill=sourceblue!8, minimum width=3.2cm, minimum height=2.8cm, align=left, inner sep=5pt},
  tamperblk/.style={draw=trapred!80, rounded corners=2pt, fill=trapred!8, minimum width=3.2cm, minimum height=2.8cm, align=left, inner sep=5pt},
  arr/.style={-{Stealth[length=2.5mm]}, very thick, quotientorange}
]
  \newcommand{\fld}[1]{{\scriptsize\ttfamily #1}}
  % Block k-1
  \node[blk] (b0) at (0,0) {
    \textbf{Block $k{-}1$ Header}\\[2pt]
    \fld{prev: $h_{k-2}$}\\
    \fld{root: $r_{k-1}$}\\
    \fld{time: $t_{k-1}$}\\
    \fld{nonce: $n_{k-1}$}\\[4pt]
    \rule{2.8cm}{0.4pt}\\[2pt]
    \textbf{ID:} $h_{k-1} = \Hash(\text{hdr})$
  };

  % Block k
  \node[blk] (b1) at (4.5,0) {
    \textbf{Block $k$ Header}\\[2pt]
    \fld{prev: $h_{k-1}$}\\
    \fld{root: $r_{k}$}\\
    \fld{time: $t_{k}$}\\
    \fld{nonce: $n_{k}$}\\[4pt]
    \rule{2.8cm}{0.4pt}\\[2pt]
    \textbf{ID:} $h_{k} = \Hash(\text{hdr})$
  };

  % Block k+1
  \node[blk] (b2) at (9.0,0) {
    \textbf{Block $k{+}1$ Header}\\[2pt]
    \fld{prev: $h_{k}$}\\
    \fld{root: $r_{k+1}$}\\
    \fld{time: $t_{k+1}$}\\
    \fld{nonce: $n_{k+1}$}\\[4pt]
    \rule{2.8cm}{0.4pt}\\[2pt]
    \textbf{ID:} $h_{k+1} = \Hash(\text{hdr})$
  };

  \draw[arr] (b0.east |- b1.170) -- node[above, font=\scriptsize, text=quotientorange!90!black]{$h_{k-1}$} (b1.west |- b1.170);
  \draw[arr] (b1.east |- b2.170) -- node[above, font=\scriptsize, text=quotientorange!90!black]{$h_{k}$} (b2.west |- b2.170);

  \node[left=0.5cm of b0, font=\footnotesize] {$\cdots$};
  \node[right=0.5cm of b2, font=\footnotesize] {$\cdots$};

  % Tamper callout below
  \node[draw=trapred!70, fill=trapred!10, rounded corners=2pt, font=\scriptsize, align=center, inner sep=4pt] at (4.5, -2.4) {
    \textbf{Tamper-Evident Cascade:}\\
    Modifying a single bit in Block $k$ changes $r_k \implies$ changes $h_k \implies$ breaks \texttt{prev} in Block $k+1$.\\
    To deceive nodes, an attacker must re-mine Block $k$, Block $k+1$, and all subsequent blocks!
  };
\end{tikzpicture}
\caption{Cryptographic block chaining. Each header commits to its transactions via a Merkle root and incorporates the exact cryptographic digest of the prior header in its \texttt{prev} field. Any retrospective alteration creates an immediate mismatch that invalidates every succeeding block.}
\label{ch36:fig-chain}
\end{figure}

\begin{whybox}
Why does chaining make the past tamper-evident? Suppose Eve wants to alter one transaction inside block $k$. That changes block $k$'s Merkle root, hence its header, hence its ID $h_k$. But block $k{+}1$ stored the \emph{old} $h_k$ in its \texttt{prev} field. For the chain to stay valid, Eve must update block $k{+}1$'s \texttt{prev} to the new $h_k$ --- which changes block $k{+}1$'s ID --- which forces an update to block $k{+}2$, and so on to the tip. By second-preimage resistance she cannot find a different block $k$ body that keeps $h_k$ the same. So a one-transaction edit cascades into re-deriving every subsequent block. With proof of work each of those re-derivations costs the puzzle's worth of computation, which is the point: integrity is cheap to \emph{check} and expensive to \emph{forge}.
\end{whybox}

Let us make this concrete with absurdly small numbers so you can do the hashing in spirit, if not literally. We will use a toy hash $h(x) = \mathrm{SHA256}(x)$ but only look at the first few hex digits. Take three blocks whose ``bodies'' are just strings.

\begin{example}[A three-block toy chain]
\label{ch36:ex-chain}
Let block $0$ be the genesis: body \texttt{"genesis"}, \texttt{prev}~$=0$. Compute $h_0=\Hash(\texttt{"0"}\concat\texttt{"genesis"})$. Block $1$ has body \texttt{"Alice->Bob 10"} and \texttt{prev}~$=h_0$, giving $h_1=\Hash(h_0\concat\texttt{"Alice->Bob 10"})$. Block $2$ has body \texttt{"Bob->Carol 4"} and \texttt{prev}~$=h_1$. Now if Eve rewrites block $1$'s body to \texttt{"Alice->Eve 10"}, recomputing gives a new $h_1'\ne h_1$. Block $2$'s stored \texttt{prev} is still the old $h_1$, so a verifier recomputing block $2$'s expected \texttt{prev} sees a mismatch and rejects the chain. Eve must also recompute $h_2$, and if there were a block $3$, that too, and so on. The implement-it box below runs exactly this.
\end{example}

\section{Merkle trees: one root for a million transactions}
\label{ch36:merkle}

A block can hold thousands of transactions. We do not want the header to grow with them, and we want a lightweight client --- say a phone wallet --- to be able to prove ``transaction $T$ is in block $k$'' while downloading only the small header, not the whole block. Both needs are met by a \emph{Merkle tree}, a binary tree of hashes named for Ralph Merkle.

\subsection{Construction}
\label{ch36:merkle-build}

Take the list of transactions $T_0,T_1,\dots,T_{m-1}$. Hash each one to get the \emph{leaves} $\ell_i = \Hash(T_i)$. Then pair the leaves and hash each pair to get the next level up: $\Hash(\ell_0\concat\ell_1)$, $\Hash(\ell_2\concat\ell_3)$, and so on. Repeat, halving the count each level, until one hash remains: the \emph{Merkle root}. If a level has an odd number of nodes, duplicate the last one (Bitcoin's convention) so it pairs with itself.

\begin{figure}[H]
\centering
\begin{tikzpicture}
  % Panel A: Merkle Tree Hierarchy
  \begin{scope}[shift={(0,0)}]
    \node[font=\small\bfseries] at (2.85, 3.8) {Panel A: Merkle Tree Aggregation};
    
    \node[draw=sourceblue!80, fill=sourceblue!15, rounded corners=2pt, font=\small, minimum width=3.2cm, minimum height=0.65cm] (rt) at (2.85, 3.0) {\textbf{Merkle Root} $h_{0123}$};
    
    \node[draw=structurepurple!80, fill=structurepurple!12, rounded corners=2pt, font=\scriptsize, minimum width=2.2cm, minimum height=0.55cm] (h01) at (1.2, 1.8) {$h_{01} = \Hash(\ell_0 \mathbin{\Vert} \ell_1)$};
    \node[draw=structurepurple!80, fill=structurepurple!12, rounded corners=2pt, font=\scriptsize, minimum width=2.2cm, minimum height=0.55cm] (h23) at (4.5, 1.8) {$h_{23} = \Hash(\ell_2 \mathbin{\Vert} \ell_3)$};
    
    \node[draw=gray!70, fill=gray!8, rounded corners=2pt, font=\tiny, minimum width=1.25cm, minimum height=0.5cm] (l0) at (0.3, 0.6) {$\ell_0 = \Hash(T_0)$};
    \node[draw=gray!70, fill=gray!8, rounded corners=2pt, font=\tiny, minimum width=1.25cm, minimum height=0.5cm] (l1) at (2.0, 0.6) {$\ell_1 = \Hash(T_1)$};
    \node[draw=gray!70, fill=gray!8, rounded corners=2pt, font=\tiny, minimum width=1.25cm, minimum height=0.5cm] (l2) at (3.7, 0.6) {$\ell_2 = \Hash(T_2)$};
    \node[draw=gray!70, fill=gray!8, rounded corners=2pt, font=\tiny, minimum width=1.25cm, minimum height=0.5cm] (l3) at (5.4, 0.6) {$\ell_3 = \Hash(T_3)$};

    \draw[-{Stealth[length=2mm]}, thick, gray] (l0) -- (h01);
    \draw[-{Stealth[length=2mm]}, thick, gray] (l1) -- (h01);
    \draw[-{Stealth[length=2mm]}, thick, gray] (l2) -- (h23);
    \draw[-{Stealth[length=2mm]}, thick, gray] (l3) -- (h23);
    \draw[-{Stealth[length=2mm]}, thick, gray] (h01) -- (rt);
    \draw[-{Stealth[length=2mm]}, thick, gray] (h23) -- (rt);

    \node[font=\tiny, fill=white, inner sep=1pt] at (2.85, -0.3) {Single 32-byte commitment commits to all leaves};
  \end{scope}

  % Panel B: SPV Audit Path for T2
  \begin{scope}[shift={(7.1, 0)}]
    \node[font=\small\bfseries] at (2.85, 3.8) {Panel B: Inclusion Audit Path for $T_2$};

    \node[draw=sourceblue!80, fill=sourceblue!15, rounded corners=2pt, font=\small, minimum width=3.2cm, minimum height=0.65cm] (rt2) at (2.85, 3.0) {\textbf{Compare:} $\text{root}' \stackrel{?}{=} h_{0123}$};

    \node[draw=quotientorange!80, fill=quotientorange!18, rounded corners=2pt, font=\scriptsize\bfseries, minimum width=2.2cm, minimum height=0.55cm] (path_h01) at (1.2, 1.8) {Sibling $h_{01}$ (Path)};
    \node[draw=targetgreen!80, fill=targetgreen!12, rounded corners=2pt, font=\scriptsize, minimum width=2.2cm, minimum height=0.55cm] (re_h23) at (4.5, 1.8) {Computed $h_{23}'$};

    \node[draw=gray!40, fill=white, rounded corners=2pt, font=\tiny, text=gray!60, minimum width=1.15cm, minimum height=0.5cm] at (0.3, 0.6) {Omitted};
    \node[draw=gray!40, fill=white, rounded corners=2pt, font=\tiny, text=gray!60, minimum width=1.15cm, minimum height=0.5cm] at (2.0, 0.6) {Omitted};
    \node[draw=targetgreen!80, fill=targetgreen!18, rounded corners=2pt, font=\scriptsize\bfseries, minimum width=1.25cm, minimum height=0.5cm] (target_l2) at (3.7, 0.6) {Leaf $\ell_2'$};
    \node[draw=quotientorange!80, fill=quotientorange!18, rounded corners=2pt, font=\scriptsize\bfseries, minimum width=1.25cm, minimum height=0.5cm] (path_l3) at (5.4, 0.6) {Sibling $\ell_3$};

    \draw[-{Stealth[length=2mm]}, very thick, targetgreen!80!black] (target_l2) -- (re_h23);
    \draw[-{Stealth[length=2mm]}, very thick, quotientorange] (path_l3) -- (re_h23);
    \draw[-{Stealth[length=2mm]}, very thick, quotientorange] (path_h01) -- (rt2);
    \draw[-{Stealth[length=2mm]}, very thick, targetgreen!80!black] (re_h23) -- (rt2);

    \node[draw=targetgreen!80, fill=targetgreen!10, rounded corners=2pt, font=\tiny, align=center, inner sep=3pt] at (2.85, -0.5) {
      \textbf{SPV Light Client Proof:}\\
      Transfers only $\{\ell_3, h_{01}\}$ ($O(\log m)$ hashes).\\
      Validates inclusion without downloading block body!
    };
  \end{scope}
\end{tikzpicture}
\caption{Merkle tree commitment and verification. Panel A depicts the full tree where $m$ transaction leaves are paired and hashed up to the 32-byte Merkle root stored in the block header. Panel B illustrates an SPV inclusion proof for transaction $T_2$: the light client receives only the audit siblings $\ell_3$ and $h_{01}$ (shaded orange), recomputes $h_{23}' = \Hash(\ell_2' \mathbin{\Vert} \ell_3)$ and $\text{root}' = \Hash(h_{01} \mathbin{\Vert} h_{23}')$, and confirms equality with the known block root.}
\label{ch36:fig-merkle}
\end{figure}

\subsection{Membership proofs}
\label{ch36:merkle-proof}

The payoff is the \emph{Merkle proof} (also called an inclusion or audit proof). To convince someone that $T_2$ is one of the leaves under a root they already trust, you do not send all the transactions. You send $T_2$ plus the \emph{sibling hashes} along the path from $T_2$'s leaf up to the root --- one hash per level. The verifier recomputes the path and checks it lands on the trusted root.

For the tree in Figure~\ref{ch36:fig-merkle}, the proof for $T_2$ is the pair $(\ell_3,\ h_{01})$ together with which side each sits on. The verifier computes
\[
  \ell_2' = \Hash(T_2),\quad
  h_{23}' = \Hash(\ell_2' \concat \ell_3),\quad
  \text{root}' = \Hash(h_{01} \concat h_{23}'),
\]
and accepts iff $\text{root}' = \text{root}$. Two hashes sent, two hashes computed, for a four-leaf tree; in general a tree of $m$ leaves needs a proof of only $\ceil{\log_2 m}$ hashes. A block with a million transactions yields a $20$-hash proof --- about $640$ bytes --- versus tens of megabytes for the full block.

\begin{whybox}
Why is the proof sound --- why can't a prover fake membership for a transaction that is not in the tree? Because forging a valid path to a fixed root requires a hash collision. Suppose Eve wants the verifier to accept a bogus $T_2^\star$. She must supply sibling hashes so that the recomputed root equals the real root. But the real root was the hash of the real $h_{01}\concat h_{23}$. To hit it with a different left-or-right input is a (second-)preimage break on the top hash; to instead make $h_{23}'$ equal the real $h_{23}$ from a different $\ell_2'$ is a collision on that node's hash. Collision resistance forbids both. So under a hash we believe is collision resistant, the only way to produce an accepting proof for $T_2^\star$ is for $T_2^\star$ to genuinely be the leaf --- the tree is a \emph{binding} commitment to its entire leaf set.
\end{whybox}

\begin{example}[A four-leaf tree, worked with small hex]
\label{ch36:ex-merkle}
We will use a 1-byte toy hash $H(x) = \mathrm{SHA256}(x)\bmod 256$, written as two hex digits, just to make the arithmetic visible (a real tree uses the full 256-bit output). Suppose the four transactions hash to leaves
\[
  \ell_0=\texttt{a3},\quad \ell_1=\texttt{7c},\quad \ell_2=\texttt{1f},\quad \ell_3=\texttt{d8}.
\]
Then $h_{01}=H(\texttt{a3}\concat\texttt{7c})$, say $=\texttt{4b}$, and $h_{23}=H(\texttt{1f}\concat\texttt{d8})$, say $=\texttt{e0}$, and the root is $H(\texttt{4b}\concat\texttt{e0})$, say $=\texttt{29}$. The proof for $T_2$ is $\{\ell_3=\texttt{d8}\ (\text{right}),\ h_{01}=\texttt{4b}\ (\text{left})\}$. A verifier holding only the root \texttt{29} and the transaction $T_2$ computes $\ell_2=H(T_2)=\texttt{1f}$, then $h_{23}'=H(\texttt{1f}\concat\texttt{d8})=\texttt{e0}$, then $\text{root}'=H(\texttt{4b}\concat\texttt{e0})=\texttt{29}$, matches, and accepts --- never having seen $T_0,T_1,T_3$.
\end{example}

\begin{applicationbox}
This is exactly how Bitcoin's \emph{Simplified Payment Verification} (SPV) clients work, the design Nakamoto called out in the whitepaper. A phone wallet stores only the $80$-byte block headers (a few dozen megabytes for the whole chain) and asks a full node for a Merkle proof that its transaction is in a block. It checks the proof against the header's root and confirms the header is buried under enough later blocks. It never downloads the multi-hundred-gigabyte full chain. Ethereum pushes the idea further with \emph{Merkle--Patricia tries} that commit not just to a list of transactions but to the entire world state --- every account balance and contract storage slot --- so a light client can be given a proof of any single account's balance. Certificate Transparency logs (Chapter 24) and the \texttt{git} version-control system are Merkle trees too: a git commit hash is a root over the whole repository tree.
\end{applicationbox}

\begin{trapbox}
Merkle trees have bitten implementers. Bitcoin's tree duplicates the last node when a level has an odd count, and this created \emph{CVE-2012-2459}: a block could be padded with a duplicate transaction so two different transaction lists produced the \emph{same} root, letting an attacker craft a valid-looking-but-invalid block that split nodes' views. The deeper lesson is that a Merkle root only commits to the \emph{multiset of leaf hashes in their tree positions}, not to a canonical list; if your encoding lets two distinct lists map to one tree, you have handed an attacker a collision the hash function never gave them. Always domain-separate leaf hashes from internal-node hashes (prefix a leaf with a \texttt{0x00} byte and an internal node with \texttt{0x01}) so a leaf can never be reinterpreted as an internal node --- a \emph{second-preimage} attack on naive Merkle trees that the prefix trick defeats.
\end{trapbox}

\begin{implementbox}
A toy chain with Merkle roots and a tamper check. \textbf{Learning only --- this omits real serialization, signatures, consensus, and uses Python's slow hashing; never run a real ledger on it.}
\begin{lstlisting}[language=Python]
import hashlib

def H(b: bytes) -> bytes:
    # SHA-256d, Bitcoin style. Domain-separate leaves vs nodes in real code.
    return hashlib.sha256(hashlib.sha256(b).digest()).digest()

def merkle_root(txs):
    level = [H(b"\x00" + t.encode()) for t in txs]   # 0x00 = leaf tag
    if not level:
        return H(b"")
    while len(level) > 1:
        if len(level) % 2:           # odd: duplicate last (Bitcoin convention)
            level.append(level[-1])
        level = [H(b"\x01" + level[i] + level[i+1])   # 0x01 = internal tag
                 for i in range(0, len(level), 2)]
    return level[0]

def block_id(prev, root, nonce):
    return H(prev + root + nonce.to_bytes(8, "big"))

# Build a 3-block chain.
chain = []
prev = bytes(32)                                   # genesis prev = 32 zero bytes
for txs in [["genesis"], ["Alice->Bob 10"], ["Bob->Carol 4", "Dave->Eve 1"]]:
    root = merkle_root(txs)
    bid  = block_id(prev, root, nonce=0)
    chain.append({"prev": prev, "root": root, "txs": txs, "id": bid})
    prev = bid

def verify(chain):
    prev = bytes(32)
    for b in chain:
        if b["prev"] != prev:                      # broken link?
            return False
        if b["root"] != merkle_root(b["txs"]):     # tampered body?
            return False
        if b["id"] != block_id(b["prev"], b["root"], 0):
            return False
        prev = b["id"]
    return True

print(verify(chain))                  # True
chain[1]["txs"] = ["Alice->Eve 10"]   # Eve rewrites history
print(verify(chain))                  # False: root no longer matches, link snaps
\end{lstlisting}
\end{implementbox}

\section{Keys, addresses, and wallets}
\label{ch36:addresses}

A coin is controlled by whoever can produce a valid signature under a particular key. So ``owning'' bitcoin means holding a private key. The chain of derivations from key to address is pure cryptography, and getting it right is what separates a wallet from a way to lose money.

\subsection{From private key to public key}
\label{ch36:keys}

Bitcoin and Ethereum both use the elliptic curve \texttt{secp256k1} (Chapter 21), defined over the prime field $\F_p$ with
\[
  p = 2^{256} - 2^{32} - 977,\qquad y^2 = x^3 + 7 \pmod p,
\]
and a fixed base point $G$ of large prime order $n \approx 2^{256}$. A private key is a random integer $d$ with $1 \le d < n$. The public key is the curve point
\[
  Q = dG,
\]
$G$ added to itself $d$ times under the curve's group law. Recovering $d$ from $Q$ is the elliptic-curve discrete log problem, believed to cost about $2^{128}$ operations --- the security floor of the whole system. The private key is $32$ bytes of entropy; lose it and the coins are gone, leak it and they are stolen. There is no reset.

\begin{example}[A microscopic curve to fix the idea]
\label{ch36:ex-curve}
\texttt{secp256k1}'s numbers are too big to trace, so work over $\F_{17}$ with $y^2 = x^3 + 7$. Checking points, $(1,5)$ is on the curve since $5^2 = 25 \equiv 8$ and $1^3+7 = 8$. Take $G=(1,5)$. Using the curve addition rules from Chapter 21, doubling gives $2G$, tripling $3G$, and so on, cycling through the group of points. If your private key is $d=4$, your public key is $Q = 4G$, computed by two doublings. An attacker who sees $Q$ but not $d$ would, on this tiny curve, just try all $d$ --- there are only a handful --- which is exactly why real curves use $n\approx 2^{256}$: the only known attack is to try (cleverly) all $d$, and there are too many.
\end{example}

\subsection{From public key to address}
\label{ch36:addr}

You do not publish $Q$ directly as your payment handle; you publish a \emph{hash} of it, called an \emph{address}. Figure~\ref{ch36:fig-address} shows the full pipeline. Bitcoin's legacy ``P2PKH'' address is
\[
  \text{address} = \mathrm{Base58Check}\big(\,\texttt{0x00} \concat \mathrm{RIPEMD160}(\mathrm{SHA256}(Q))\,\big),
\]
that is, SHA-256 then RIPEMD-160 (jointly nicknamed HASH160) to get a $20$-byte fingerprint, prefixed with a version byte, then Base58Check-encoded with a $4$-byte checksum so typos do not send coins into the void. Ethereum addresses are the last $20$ bytes of $\mathrm{Keccak256}(Q)$.

\begin{figure}[H]
\centering
\begin{tikzpicture}[font=\footnotesize,
  box/.style={draw=sourceblue!80, fill=sourceblue!12, rounded corners=2pt, minimum height=8mm, align=center, inner sep=4pt},
  op/.style={draw=structurepurple!80, fill=structurepurple!12, rounded corners=2pt, minimum height=6mm, inner sep=3pt, align=center},
  res/.style={draw=targetgreen!80, fill=targetgreen!12, rounded corners=2pt, minimum height=8mm, align=center, inner sep=4pt},
  ar/.style={-{Stealth[length=2.2mm]}, thick}]
  \node[box] (d) at (0, 1.5) {\textbf{Private Key} $d$\\$256$ bits entropy};
  \node[op] (mul) at (4.5, 1.5) {$\times G$ on \texttt{secp256k1}\\Discrete Log Trapdoor};
  \node[box] (q) at (9.2, 1.5) {\textbf{Public Key} $Q = dG$\\$33$ bytes (compressed)};
  \node[op] (h160) at (9.2, 0.0) {$\Hash_{160}(Q) = \mathrm{RIPEMD160}\bigl(\mathrm{SHA256}(Q)\bigr)$\\Preimage Resistance Trapdoor};
  \node[box] (fp) at (4.5, 0.0) {\textbf{Key Hash}\\$20$ bytes};
  \node[op] (enc) at (4.5, -1.5) {Version tag \texttt{0x00} + $4$-byte Checksum\\Base58Check};
  \node[res] (addr) at (0, -1.5) {\textbf{Public Address}\\e.g., \texttt{1A1zP1...}};
  
  \draw[ar, sourceblue!80!black] (d) -- (mul);
  \draw[ar, sourceblue!80!black] (mul) -- (q);
  \draw[ar, structurepurple!80!black] (q) -- (h160);
  \draw[ar, structurepurple!80!black] (h160) -- (fp);
  \draw[ar, structurepurple!80!black] (fp) -- (enc);
  \draw[ar, targetgreen!80!black] (enc) -- (addr);
\end{tikzpicture}
\caption{The one-way derivation pipeline from private key to payment address. Each operational stage is mathematically irreversible: curve multiplication conceals $d$ under the elliptic-curve discrete log problem, and the dual-hash ($\Hash_{160}$) pipeline conceals $Q$ behind preimage resistance until the coins are spent.}
\label{ch36:fig-address}
\end{figure}

\begin{whybox}
Why hash the public key instead of using it directly? Three reasons. (1) \emph{Brevity}: $20$ bytes is shorter than a $33$-byte compressed public key, friendlier to type and to print as a QR code. (2) \emph{A second wall}: if elliptic-curve discrete log were ever broken, an address whose public key has never been revealed still hides $Q$ behind a preimage-resistant hash --- the public key only becomes visible \emph{when you spend}, so an unspent, never-reused address has a hash-based backstop. This is the practical case against address reuse. (3) \emph{Checksumming}: the Base58Check wrapper carries a checksum that catches mistyped addresses before funds move. Note the security floor for the hash wall is $2^{80}$ (RIPEMD-160 is $160$ bits, collisions at $2^{80}$), weaker than the curve's $2^{128}$; this is a known and debated soft spot, mitigated in practice by not reusing addresses.
\end{whybox}

\subsection{Wallets and key derivation}
\label{ch36:wallets}

A \emph{wallet} is a key manager, not a place where coins live --- the coins are entries on the chain; the wallet holds the keys that authorize spending them. Managing one random key per address is unwieldy, so modern wallets are \emph{hierarchical deterministic} (HD, BIP-32): a single master seed deterministically derives an unlimited tree of keys via repeated HMAC-SHA512. The seed itself is usually encoded as a \emph{mnemonic} (BIP-39) --- the famous twelve or twenty-four English words --- whose entropy plus a checksum maps to the seed.

\begin{intuitionbox}
The whole hierarchy --- millions of addresses across many coins --- collapses back to twelve words. That compactness is both its strength and its hazard. Back up the words once and you can regenerate every key forever, on any compatible wallet, even if your device is destroyed. But anyone who reads those words owns everything the seed controls, instantly and irreversibly. The mnemonic is not a password you can change; it is the root of a deterministic tree, so ``rotating'' it means moving funds to a new tree. Treat the twelve words the way you would treat the only key to a vault with no locksmith.
\end{intuitionbox}

\begin{applicationbox}
HD derivation has a subtle but important split: \emph{hardened} vs \emph{non-hardened} child keys. Non-hardened derivation lets you compute child \emph{public} keys from a parent public key without the private key --- useful for a webshop that wants a fresh receiving address per customer without keeping any private key on the public server. The catch: if an attacker learns one non-hardened child \emph{private} key and the parent \emph{public} key, the BIP-32 math lets them recover the parent private key and thus every sibling. Hardened derivation breaks that link by feeding the parent private key into the HMAC, at the cost of not supporting public-only derivation. Real wallets use hardened derivation for the account-level nodes and non-hardened only for the leaf addresses under an account, so a leaked leaf key cannot climb the tree.
\end{applicationbox}

\section{Signatures: authorizing a spend}
\label{ch36:signatures}

A transaction is an explicit authorization: ``consume these specific unspent outputs that I control, and create new unspent outputs assigned to these recipients.'' In the Unspent Transaction Output (UTXO) model, coins do not exist as account balances; they exist as atomic, discrete chunks of value called UTXOs. To spend a UTXO, the transaction must satisfy the locking script associated with it, which almost universally requires presenting a valid cryptographic signature under the private key matching the locked address.

Figure~\ref{ch36:fig-utxo} illustrates this ledger mechanism. Panel A traces a legitimate spend consuming an existing UTXO, while Panel B demonstrates why double-spending fails immediately at the peer-to-peer relay layer.

\begin{figure}[H]
\centering
\begin{tikzpicture}
  % Panel A: Legitimate UTXO spend
  \begin{scope}[shift={(0,0)}]
    \node[font=\small\bfseries] at (3.4, 3.6) {Panel A: Valid State Transition (UTXO Spend)};
    
    % Input UTXO
    \node[draw=sourceblue!80, fill=sourceblue!12, rounded corners=2pt, font=\scriptsize, align=center, inner sep=4pt] (utxo_in) at (-1.5, 2.0) {
      \textbf{Existing UTXO}\\
      Value: $10.0$ BTC\\
      Locked to: $\Hash(Q_{\text{Alice}})$
    };

    % Tx Box
    \node[draw=structurepurple!80, fill=structurepurple!10, rounded corners=2pt, font=\scriptsize, align=center, inner sep=5pt] (tx) at (3.2, 2.0) {
      \textbf{Transaction $\text{Tx}_1$}\\
      Input: $\text{UTXO}_{\text{ref}}$\\
      Witness: $\sigma_{\text{Alice}}$ on $\Hash(\text{Tx}_1)$
    };

    % Outputs
    \node[draw=targetgreen!80, fill=targetgreen!12, rounded corners=2pt, font=\scriptsize, align=left, inner sep=4pt] (out1) at (8.0, 2.65) {
      \textbf{New UTXO 1:} $7.0$ BTC\\
      Locked to $\Hash(Q_{\text{Bob}})$
    };
    \node[draw=targetgreen!80, fill=targetgreen!12, rounded corners=2pt, font=\scriptsize, align=left, inner sep=4pt] (out2) at (8.0, 1.35) {
      \textbf{New UTXO 2:} $2.99$ BTC\\
      Change $\to \Hash(Q_{\text{Alice}})$
    };

    \draw[-{Stealth[length=2mm]}, thick] (utxo_in.east) -- (tx.west);
    \draw[-{Stealth[length=2mm]}, thick] (tx.east) -- (out1.west);
    \draw[-{Stealth[length=2mm]}, thick] (tx.east) -- (out2.west);

    \node[draw=targetgreen!80, fill=targetgreen!8, rounded corners=2pt, font=\scriptsize, align=center, inner sep=3.5pt] at (3.4, 0.2) {
      $\text{UTXO}_{\text{Alice}}$ is consumed and pruned from UTXO set. Miner collects implicit fee: $10.0 - (7.0 + 2.99) = 0.01$ BTC.
    };
  \end{scope}

  % Panel B: Attempted Double-Spend (stacked below Panel A)
  \begin{scope}[shift={(0, -2.5)}]
    \node[font=\small\bfseries] at (3.4, 1.8) {Panel B: Double-Spend Rejection};

    \node[draw=trapred!80, fill=trapred!10, rounded corners=2pt, font=\scriptsize, align=center, inner sep=5pt] (tx_bad) at (-0.4, 0.7) {
      \textbf{Compromised $\text{Tx}_2$ (Eve)}\\
      Attempts to spend $\text{UTXO}_{\text{Alice}}$ a second time!\\
      Valid signature, but references spent coin.
    };

    \node[draw=trapred!80, fill=trapred!15, rounded corners=2pt, font=\scriptsize\bfseries, align=center, inner sep=5pt] (verdict) at (7.4, 0.7) {
      \textbf{REJECTED BY EVERY NODE}\\[2pt]
      Lookup in global UTXO database shows\\
      $\text{UTXO}_{\text{Alice}}$ has already been spent.
    };

    \draw[-{Stealth[length=2.2mm]}, very thick, trapred] (tx_bad.east) -- node[above, font=\tiny\bfseries, trapred]{UTXO missing} (verdict.west);
  \end{scope}
\end{tikzpicture}
\caption{The Unspent Transaction Output (UTXO) accounting model and double-spending prevention. In Panel A, an unspent output is consumed in its entirety, validated by Alice's signature, and transmuted into new unspent outputs for Bob and Alice's change address. In Panel B, any concurrent or subsequent attempt to spend the same referenced UTXO fails consensus validation because full nodes track the global unspent set and reject previously spent references.}
\label{ch36:fig-utxo}
\end{figure}

\subsection{ECDSA}
\label{ch36:ecdsa}

Bitcoin's original signature scheme is ECDSA. With private key $d$, public key $Q=dG$, group order $n$, and message hash $z = \Hash(\text{transaction})$ reduced mod $n$, signing produces a pair $(r,s)$:

\begin{enumerate}
  \item Pick a random per-signature \emph{nonce} $k$, with $1 \le k < n$.
  \item Compute the point $kG = (x_1,y_1)$ and set $r = x_1 \bmod n$. If $r=0$, restart.
  \item Compute $s = k^{-1}(z + r\,d) \bmod n$. If $s=0$, restart.
  \item The signature is $(r,s)$.
\end{enumerate}

Verification, given $(r,s)$, $z$, and $Q$: compute $u_1 = z\,s^{-1}\bmod n$, $u_2 = r\,s^{-1}\bmod n$, form the point $u_1 G + u_2 Q = (x_1,y_1)$, and accept iff $x_1 \bmod n = r$.

\begin{whybox}
The verification identity falls out of the algebra. From $s = k^{-1}(z+rd)$ we get $k = s^{-1}(z+rd) = s^{-1}z + s^{-1}rd = u_1 + u_2 d \pmod n$. Multiply both sides by $G$: $kG = u_1 G + u_2 (dG) = u_1 G + u_2 Q$. So the point the verifier reconstructs from public data, $u_1 G + u_2 Q$, equals the point $kG$ the signer used, and its $x$-coordinate mod $n$ is exactly $r$. Only someone who knows $d$ can produce an $s$ that makes this close, because $s$ binds $z$, $r$, and $d$ together through the secret $k$; forging $(r,s)$ for a chosen $z$ without $d$ requires solving the elliptic-curve discrete log.
\end{whybox}

\begin{attackbox}
\textbf{Nonce reuse is fatal.} The per-signature $k$ must be \emph{secret} and \emph{never reused}. Reuse $k$ across two signatures $(r,s_1)$ and $(r,s_2)$ on messages $z_1\ne z_2$ --- note the $r$ is identical, a dead giveaway (Figure~\ref{ch36:fig-nonce-reuse}) --- and the private key falls out by hand:
\[
  s_1 - s_2 = k^{-1}(z_1 - z_2) \;\Rightarrow\; k = \frac{z_1 - z_2}{s_1 - s_2} \bmod n,
  \qquad d = \frac{s_1 k - z_1}{r}\bmod n.
\]
This is not theoretical. In 2010 the Sony PlayStation~3 used a \emph{fixed} $k$ for its ECDSA code-signing key; the master key was extracted and the console's chain of trust collapsed. In 2013 a flaw in Android's \texttt{SecureRandom} produced repeated $k$ values in some wallet apps and real bitcoin was stolen from the affected addresses. The fix is \emph{deterministic nonces} (RFC~6979): derive $k$ via HMAC from the key $d$ and message $z$ so it is reproducible, unique per message, and never depends on a possibly-broken RNG. Modern Bitcoin signing uses RFC~6979 universally.
\end{attackbox}

\begin{figure}[H]
\centering
\begin{tikzpicture}[font=\footnotesize,
  sig/.style={draw=sourceblue!80, fill=sourceblue!8, rounded corners=2pt, minimum width=4.0cm, minimum height=1.1cm, align=center},
  stp/.style={draw=trapred!80, fill=trapred!10, rounded corners=2pt, minimum width=3.6cm, minimum height=1.1cm, align=center},
  arr/.style={-{Stealth[length=2.2mm]}, thick, trapred!85!black}
]
  \node[sig] (s1) at (0, 0.75) {\textbf{Signature 1} (same $k, r$)\\$s_1 = k^{-1}(z_1 + rd) \bmod n$};
  \node[sig] (s2) at (0, -0.75) {\textbf{Signature 2} (same $k, r$)\\$s_2 = k^{-1}(z_2 + rd) \bmod n$};

  \node[stp] (rec_k) at (4.9, 0) {\textbf{1. Subtract ($rd$ cancels)}\\$k = (z_1 - z_2)(s_1 - s_2)^{-1}$};
  \node[stp, draw=trapred, fill=trapred!18] (rec_d) at (9.5, 0) {\textbf{2. Solve for Secret Key}\\$d = r^{-1}(s_1 k - z_1) \bmod n$};

  \draw[arr] (s1.east) -- (rec_k.165);
  \draw[arr] (s2.east) -- (rec_k.195);
  \draw[arr] (rec_k.east) -- (rec_d.west);
\end{tikzpicture}
\caption{Why ECDSA nonce reuse is fatal. Two signatures generated with the same nonce $k$ share the same $r$ and give two linear equations in two unknowns ($k$ and $d$): subtracting eliminates $rd$ to expose $k$, which immediately unlocks the long-term private key $d$.}
\label{ch36:fig-nonce-reuse}
\end{figure}

\begin{example}[ECDSA on a toy curve]
\label{ch36:ex-ecdsa}
Take a curve with group order $n=19$ and a base point $G$, private key $d=7$. Suppose for message hash $z=11$ the signer's RNG yields $k=3$, and $kG$ has $x$-coordinate $x_1=10$, so $r = 10 \bmod 19 = 10$. Then $s = k^{-1}(z+rd) \bmod 19$. Compute $k^{-1} = 3^{-1} \bmod 19 = 13$ (since $3\cdot13=39\equiv1$). And $z+rd = 11 + 10\cdot7 = 81 \equiv 81 - 4\cdot19 = 5 \pmod{19}$. So $s = 13\cdot5 = 65 \equiv 65 - 3\cdot19 = 8 \pmod{19}$. Signature $(r,s)=(10,8)$. A verifier computes $s^{-1}=8^{-1}\bmod 19 = 12$ (since $8\cdot12=96\equiv1$), then $u_1 = z s^{-1} = 11\cdot12 = 132 \equiv 132-6\cdot19=18$, $u_2 = r s^{-1} = 10\cdot12 = 120 \equiv 120-6\cdot19=6$, forms $18G + 6Q$, and checks its $x$-coordinate mod $19$ equals $10$. It does, because $u_1 + u_2 d = 18 + 6\cdot7 = 60 \equiv 60-3\cdot19 = 3 = k \pmod{19}$, so $18G+6Q = 3G = kG$, whose $x$-coordinate is the $10$ we started with.
\end{example}

\subsection{Schnorr signatures}
\label{ch36:schnorr}

In 2021 Bitcoin's Taproot upgrade added Schnorr signatures (BIP-340) alongside ECDSA. Schnorr is older, simpler, and has a clean security proof under the discrete-log assumption (ECDSA's security proof is messier). For private key $d$, public key $Q=dG$, message $m$:

\begin{enumerate}
  \item Pick nonce $k$, compute $R = kG$.
  \item Compute challenge $e = \Hash(R \concat Q \concat m) \bmod n$.
  \item Set $s = k + e\,d \bmod n$. Signature is $(R, s)$.
\end{enumerate}

Verify by checking $sG \stackrel{?}{=} R + eQ$. The check holds because $sG = (k+ed)G = kG + e(dG) = R + eQ$. Schnorr's headline feature is \emph{linearity}: signatures add. Two signers with keys $d_1,d_2$ can produce a single Schnorr signature valid under the combined key $Q_1+Q_2$ (this is \emph{MuSig}, done carefully to resist a rogue-key attack). On-chain, a $3$-of-$3$ multisig or a Lightning channel close looks like one ordinary signature --- smaller, cheaper, and more private, since you cannot tell a multisig from a single-sig spend.

\begin{focusbox}
Both schemes prove the same thing: \emph{the signer knew the private key $d$ for this public key, and authorized this exact message $m$}. Change one satoshi of the transaction and $z$ (or $e$) changes, so the old signature no longer verifies --- signatures bind to the precise bytes. What signatures do \emph{not} prove is that the spend is wise, legal, or what the signer intended; a wallet that displays the wrong amount, or malware that swaps the destination address before signing, gets a perfectly valid signature on a transaction that robs you. The cryptography authenticates bytes, not intentions.
\end{focusbox}

\section{Proof of work: making history expensive to rewrite}
\label{ch36:pow}

Signatures and hash chaining make tampering \emph{detectable}. They do not by themselves stop an attacker from broadcasting an entirely alternative chain and claiming it is the real one. With no referee, which chain is canonical? Bitcoin's answer: the chain with the most cumulative \emph{proof of work}. Producing each block requires solving a hash puzzle that has no shortcut, so the longest valid chain is the one the most computation went into, and rewriting old blocks means out-computing the entire rest of the network.

\subsection{The puzzle}
\label{ch36:puzzle}

A block is valid only if its header hash, read as a $256$-bit number, is \emph{below a target} $\mathcal{T}$:
\[
  \Hash(\text{header}) < \mathcal{T}.
\]
The header includes the nonce, so miners try nonce after nonce (and tweak other malleable fields when the nonce space is exhausted) until the hash comes out small enough (Figure~\ref{ch36:fig-pow-target}). Because the hash is preimage-resistant, there is no way to compute a winning nonce directly; you can only guess and check. Lowering $\mathcal{T}$ makes wins rarer and the puzzle harder; the network adjusts $\mathcal{T}$ every $2016$ blocks so blocks arrive about every ten minutes on average regardless of total mining power.

\begin{figure}[H]
\centering
\begin{tikzpicture}[font=\footnotesize]
  % Full 256-bit space bar
  \fill[gray!12] (0, 0) rectangle (10.5, 0.9);
  \fill[targetgreen!25] (0, 0) rectangle (2.2, 0.9);
  \draw[thick, gray!70] (0, 0) rectangle (10.5, 0.9);
  \draw[very thick, targetgreen!80!black] (2.2, -0.15) -- (2.2, 1.05);

  \node[below, font=\scriptsize] at (0, -0.15) {$0$};
  \node[below, font=\scriptsize\bfseries, targetgreen!80!black] at (2.2, -0.15) {Target $\mathcal{T}$};
  \node[below, font=\scriptsize] at (10.5, -0.15) {$2^{256}-1$};

  \node[font=\scriptsize\bfseries, targetgreen!75!black] at (1.1, 0.45) {Valid ($<\mathcal{T}$)};
  \node[font=\scriptsize, gray!70!black] at (6.3, 0.45) {Rejected ($\Hash(\text{hdr}) \ge \mathcal{T}$)};

  % Nonce trials above
  \node[draw=trapred!75, fill=trapred!10, rounded corners=2pt, inner sep=2.5pt, font=\scriptsize] (n0) at (4.8, 1.9) {nonce $0 \mapsto \texttt{9f2a...}$};
  \node[draw=trapred!75, fill=trapred!10, rounded corners=2pt, inner sep=2.5pt, font=\scriptsize] (n1) at (8.4, 1.9) {nonce $1 \mapsto \texttt{e41c...}$};
  \node[draw=targetgreen!80!black, fill=targetgreen!15, rounded corners=2pt, inner sep=2.5pt, font=\scriptsize\bfseries] (nw) at (1.1, 1.9) {nonce $n^\star \mapsto \texttt{00003b...}$};

  \draw[-{Stealth[length=2mm]}, thick, trapred!80] (n0.south) -- (4.8, 0.92);
  \draw[-{Stealth[length=2mm]}, thick, trapred!80] (n1.south) -- (8.4, 0.92);
  \draw[-{Stealth[length=2mm]}, very thick, targetgreen!80!black] (nw.south) -- (1.1, 0.92);
\end{tikzpicture}
\caption{Proof-of-work as a target-threshold search in $\{0,\dots,2^{256}-1\}$. Each candidate nonce maps to a pseudorandom point along the $256$-bit interval; only hashes landing in the green zone $[0,\mathcal{T})$ (equivalently, starting with enough leading zero bits) produce a valid block.}
\label{ch36:fig-pow-target}
\end{figure}

\begin{intuitionbox}
``Below the target'' is the same as ``the hash starts with enough zero bits.'' If the target is $2^{256}/D$, then a random hash lands below it with probability $1/D$, so you expect to try about $D$ nonces per block. Each try is one hash; finding a hash with the first $20$ bits zero takes about $2^{20}\approx10^6$ tries, the first $76$ bits zero (a recent real Bitcoin difficulty) takes about $2^{76}$. There is no cleverness available --- preimage resistance guarantees the best strategy is brute force --- which is precisely why the work is a meaningful, unforgeable measure of effort. You cannot fake having done $2^{76}$ hashes; the winning nonce is the receipt.
\end{intuitionbox}

\begin{connectionbox}
The nonce search is a Bernoulli process: each hash is an independent trial that succeeds with probability $q = \mathcal{T}/2^{256}$. The number of tries until the first success is geometric with mean $1/q$, and over many blocks the inter-block times are (to good approximation) exponential --- a Poisson process of block discoveries. That is the probability companion's machinery applied directly: difficulty retargeting is just steering the rate parameter of a Poisson process toward one event per ten minutes. The same memorylessness explains why ``this block is taking a while'' tells you nothing about when the next one arrives.
\end{connectionbox}

\begin{example}[Mining a toy block]
\label{ch36:ex-mine}
Set an easy target: ``hash must start with four zero hex digits'' ($\mathcal{T}=2^{256}/2^{16}$, so $q=1/65536$). The header is the string \texttt{"prev|root|"} followed by a nonce. Try nonce $0,1,2,\dots$, hashing \texttt{"...|0"}, \texttt{"...|1"}, until SHA-256d of the header begins \texttt{0000}. You expect about $65{,}536$ tries; on a laptop this finishes in a fraction of a second, and you can literally watch the nonce climb. Real Bitcoin uses the same loop with $\mathcal{T}$ vastly smaller, which is why it takes the combined planet's worth of ASICs ten minutes instead of your laptop a millisecond.
\end{example}

\begin{implementbox}
Mining the toy block from Example~\ref{ch36:ex-mine}. \textbf{Learning only --- a real miner is optimized ASIC firmware; this is to feel the search.}
\begin{lstlisting}[language=Python]
import hashlib

def mine(prev_hex, root_hex, leading_zeros=2):
    prefix = "0" * leading_zeros
    nonce = 0
    while True:
        header = f"{prev_hex}|{root_hex}|{nonce}".encode()
        h = hashlib.sha256(hashlib.sha256(header).digest()).hexdigest()
        if h.startswith(prefix):
            return nonce, h
        nonce += 1

nonce, h = mine("00ab", "4f12", leading_zeros=4)   # ~65536 expected tries
print("nonce:", nonce, "hash:", h)                  # hash begins 0000...
# Verifying is one hash; finding the nonce was ~65536. That asymmetry IS the point.
\end{lstlisting}
\end{implementbox}

\subsection{Why rewriting is hard: the longest-chain rule}
\label{ch36:longest}

Rule-following miners always extend the chain with the most work. To rewrite block $k$, an attacker must redo block $k$'s proof of work \emph{and} every block after it, then catch up to and overtake the protocol-following tip --- which keeps growing while the attacker computes. If the attacker controls a fraction $\beta$ of total hash power and the rest of the network $1-\beta$, the attacker's private chain falls behind in expectation whenever $\beta < 1/2$. The gap to overcome grows with each confirmation, so the probability of a successful rewrite drops roughly geometrically in the number of blocks built on top --- which is why merchants ``wait for six confirmations.''

\begin{attackbox}
\textbf{The 51\% attack.} If $\beta > 1/2$, the attacker can on average out-produce the rest of the network, build a longer alternative chain, and replace the public history --- enabling \emph{double-spends}: pay a merchant on the public chain, take the goods, then publish a longer chain in which that payment never happened. What it does \emph{not} grant is the ability to forge signatures, mint coins out of thin air, or spend others' addresses --- those are protected by ECDSA and the issuance rules, which a longer chain cannot violate and still be valid. A 51\% attacker can reorder, censor, and reverse \emph{recent} transactions involving their own coins, not break cryptography. Smaller chains with little total hash power have suffered real 51\% attacks (Ethereum Classic, Bitcoin Gold); Bitcoin's hash rate makes it prohibitively expensive, though that is an economic guarantee, not a mathematical one. Proof-of-stake chains (Ethereum since 2022) replace hash power with staked capital and slashing penalties, shifting the same ``majority can rewrite'' threshold from compute to coins.
\end{attackbox}

\begin{figure}[H]
\centering
\begin{tikzpicture}[font=\footnotesize,
  b/.style={draw,rounded corners,minimum width=1.1cm,minimum height=0.7cm},
  hon/.style={b,fill=green!8},att/.style={b,fill=red!8}]
  \node[hon] (h0) {$k{-}1$};
  \node[hon,right=0.5cm of h0] (h1) {$k$};
  \node[hon,right=0.5cm of h1] (h2) {$k{+}1$};
  \node[hon,right=0.5cm of h2] (h3) {$k{+}2$};
  \node[hon,right=0.5cm of h3] (h4) {$k{+}3$};
  \draw[-{Stealth[length=2mm]}] (h0)--(h1); \draw[-{Stealth[length=2mm]}] (h1)--(h2);
  \draw[-{Stealth[length=2mm]}] (h2)--(h3); \draw[-{Stealth[length=2mm]}] (h3)--(h4);
  \node[att,below=0.7cm of h1] (a1) {$k'$};
  \node[att,right=0.5cm of a1] (a2) {$k{+}1'$};
  \draw[-{Stealth[length=2mm]}] (h0)|-(a1); \draw[-{Stealth[length=2mm]}] (a1)--(a2);
  \node[right=0.6cm of a2,red!60!black] {attacker's private fork (behind)};
  \node[right=0.3cm of h4,green!50!black] {protocol-following tip (ahead)};
\end{tikzpicture}
\caption{An attacker forking from block $k{-}1$ must outrun the protocol-following chain. While the attacker built two blocks, the rest of the network built four. With under half the hash power, the attacker's fork falls further behind every block, so deeply-buried history is effectively permanent.}
\label{ch36:fig-fork}
\end{figure}

\section{What the cryptography does and does not guarantee}
\label{ch36:guarantees}

This is the section to read twice. Blockchain marketing routinely claims properties the math never promised. Be precise about the boundary.

\textbf{What the cryptography \emph{does} guarantee} (assuming SHA-256, RIPEMD-160, and \texttt{secp256k1} hold, and a protocol-following majority for consensus):

\begin{itemize}
  \item \emph{Integrity of recorded history.} Hash chaining plus proof of work make editing a buried transaction infeasible. The ledger is tamper-evident, and tamper-resistant in proportion to the work piled on top.
  \item \emph{Authorization.} Only the holder of a private key can produce a spending signature for its address. Coins cannot be stolen without the key (or a break in ECDSA).
  \item \emph{Public verifiability.} Anyone can re-derive every hash, check every signature, and confirm every Merkle proof from scratch, trusting no participant. Correctness is checkable, not asserted.
  \item \emph{Commitment / non-repudiation of the bytes.} A signed transaction is a binding statement: the signer cannot later deny authorizing those exact bytes.
\end{itemize}

\textbf{What it does \emph{not} guarantee:}

\begin{itemize}
  \item \emph{Value.} No theorem says a coin is worth anything. Cryptography secures the ledger's \emph{integrity}, not the \emph{price} or liquidity of what it records. A perfectly secured ledger of worthless tokens is still worthless.
  \item \emph{Correctness of intent.} Signatures bind bytes, not meaning. Malware that rewrites the recipient address before you sign, or a smart contract with a logic bug, produces a cryptographically flawless theft. The 2016 DAO hack drained tens of millions of dollars through a contract reentrancy bug --- the signatures and hashes were all valid; the \emph{logic} was wrong.
  \item \emph{Privacy.} The default ledger is public and pseudonymous, not anonymous. Addresses are not names, but transaction-graph analysis routinely de-anonymizes users. Confidentiality requires extra machinery (zero-knowledge proofs, Chapter 31; mixing; confidential transactions).
  \item \emph{Key safety.} Lose the private key (or the BIP-39 mnemonic) and the coins are unrecoverable --- there is no password reset and no support line. Leak it and they are stolen. Millions of bitcoin are estimated to be permanently lost to forgotten keys. The cryptography has no opinion about your backups.
  \item \emph{Legality or off-chain truth.} The chain can record a transaction that is fraud, theft, or fiction. It guarantees the \emph{record} is consistent, not that the recorded claim is true or lawful. ``Code is law'' is a slogan, not a fact.
  \item \emph{Liveness against a hash-power majority, forever.} The protocol-following-majority assumption is economic, not mathematical. A 51\% attacker breaks no cipher but can reorder and reverse recent history. Security here is ``expensive to attack,'' not ``impossible.''
\end{itemize}

\begin{trapbox}
The most expensive misunderstanding in the space is conflating ``the cryptography is unbreakable'' with ``my money is safe.'' The cryptography being sound is necessary, not sufficient. Real losses overwhelmingly come from the \emph{edges} the math does not touch: phishing the seed phrase, malware swapping the destination address at signing time, exchange custodians getting hacked, smart-contract bugs, lost backups, and social-engineering a wallet's recovery flow. A cryptographer's job on a blockchain team is as much about these seams --- key management, signing-device UX, contract auditing --- as about the curve arithmetic.
\end{trapbox}

\begin{applicationbox}
Where these primitives show up beyond cryptocurrency. \emph{Certificate Transparency} (Chapter 24) is a Merkle-tree append-only log of every TLS certificate issued, letting anyone detect a rogue certificate; it borrows the blockchain's tamper-evident-log idea without the coin or the proof of work. \emph{Git} chains commits by hash exactly as a blockchain chains blocks --- a commit hash is a root over the tree plus the parent commit hash, which is why rewriting git history changes every later commit ID. Supply-chain and software-provenance systems (Sigstore, in-toto) use Merkle logs to bind artifacts to their build records. The reusable lesson: \emph{a hash-chained Merkle log gives you a tamper-evident, publicly verifiable history.} Whether you bolt on a coin and a proof-of-work consensus depends on whether you have a trusted operator (then you do not need them) or a permissionless crowd (then you do).
\end{applicationbox}

\section{Putting it together: the life of a transaction}
\label{ch36:lifecycle}

Trace one payment through every primitive in this chapter.

\begin{enumerate}
  \item \textbf{Keys.} Alice's wallet holds a private key $d$; her address is $\mathrm{HASH160}(dG)$, Base58Check-encoded. Bob gives Alice his address.
  \item \textbf{Build.} Alice's wallet constructs a transaction: inputs referencing coins she controls, outputs paying Bob's address (and change back to herself).
  \item \textbf{Sign.} The wallet hashes the transaction to $z$ and produces an ECDSA (or Schnorr) signature with $d$, using an RFC~6979 deterministic nonce. The signature proves authorization without revealing $d$.
  \item \textbf{Broadcast.} The signed transaction floods the peer-to-peer network. Every node checks the signature against the input public keys and confirms the coins are unspent.
  \item \textbf{Merkle.} A miner collects valid transactions into a block, hashes each, and builds the Merkle tree; the root goes in the header.
  \item \textbf{Chain \& mine.} The header includes the previous block's hash (chaining) and a nonce; the miner searches nonces until $\Hash(\text{header}) < \mathcal{T}$ (proof of work).
  \item \textbf{Confirm.} The block propagates; nodes verify the proof of work, the Merkle root, and every signature, then extend their chain. Each subsequent block buries Alice's payment deeper, making reversal exponentially harder.
  \item \textbf{Verify later.} Bob's phone, an SPV client, holds only headers. It asks for a Merkle proof that Alice's transaction is in the block and checks it against the header's root --- never downloading the full chain.
\end{enumerate}

Six primitives, one payment: hash functions glue and commit, Merkle trees compress and prove membership, elliptic-curve keys identify, signatures authorize, hash chaining makes the past tamper-evident, and proof of work makes rewriting it expensive. Remove any one and the structure fails: no signatures and anyone spends your coins; no chaining and the past is editable; no proof of work and any node rewrites history for free; no Merkle trees and light clients must trust full nodes. The blockchain is the interlock.

Correct algorithms on paper still have to run on physical hardware. In Part~VI we turn to what happens when they do: Chapter~37 opens with side-channel and timing attacks, showing how execution time, cache hits, and power draw can leak a secret key without touching the math.

\section*{Exercises}
\addcontentsline{toc}{section}{Exercises}

\begin{exercise}[Chaining cascade]
\label{ch36:ex-cascade}
A chain has blocks $0,1,2,3,4$, each header containing \texttt{prev} (the previous header's hash) and a body. Eve edits the body of block $2$. List exactly which block IDs change and which \texttt{prev} fields a verifier finds inconsistent. If each block required $W$ units of proof-of-work, how much work must Eve redo to present a self-consistent altered chain of the same length?
\end{exercise}

\begin{exercise}[Merkle proof by hand]
\label{ch36:ex-merkleproof}
Using the toy 1-byte hash of Example~\ref{ch36:ex-merkle} with leaves $\ell_0=\texttt{a3},\ell_1=\texttt{7c},\ell_2=\texttt{1f},\ell_3=\texttt{d8}$ and the intermediate values given there, write down the Merkle proof for $T_1$ (not $T_2$). State which sibling goes on which side, and show the three-step recomputation a verifier performs to reach the root \texttt{29}.
\end{exercise}

\begin{exercise}[Odd leaves]
\label{ch36:ex-oddleaves}
A block has five transactions. Draw the Merkle tree using Bitcoin's ``duplicate the last node'' rule for odd levels. How many hashes are in the inclusion proof for the third transaction? Then explain in one or two sentences how the duplication rule enabled CVE-2012-2459, and how a leaf/internal domain-separation tag prevents it.
\end{exercise}

\begin{exercise}[Address derivation]
\label{ch36:ex-address}
Given a compressed public key as $33$ bytes, write the sequence of operations (which hash, in which order, what prefix byte, what encoding) to produce a Bitcoin P2PKH address. Why is the public key hashed rather than published directly? Give two distinct reasons, and state the bit-security of the RIPEMD-160 wall against collisions.
\end{exercise}

\begin{exercise}[Break ECDSA nonce reuse]
\label{ch36:ex-noncereuse}
Two ECDSA signatures share the same $r$: $(r,s_1)$ on $z_1$ and $(r,s_2)$ on $z_2$, with group order $n$. Derive $k$ and then the private key $d$ in terms of $r, s_1, s_2, z_1, z_2$. Then plug in $n=19$, $r=10$, $z_1=11$, $s_1=8$, $z_2=4$, $s_2=12$ and recover $d$ as a concrete number. (Check your answer against the signing equation.)
\end{exercise}

\begin{exercise}[Verify a toy ECDSA signature]
\label{ch36:ex-verify}
With $n=19$, suppose a signature is $(r,s)=(10,8)$ on $z=11$, and you are told $u_1 + u_2 d \equiv 3 \pmod{19}$ where $d$ is the signer's key. Compute $u_1$ and $u_2$ from $(r,s,z)$, and confirm the verification point reconstructs $kG$ with $k=3$. (You may reuse the inverses from Example~\ref{ch36:ex-ecdsa}.)
\end{exercise}

\begin{exercise}[Schnorr verification]
\label{ch36:ex-schnorr}
For Schnorr with private key $d$, $Q=dG$, nonce $k$, $R=kG$, challenge $e=\Hash(R\concat Q\concat m)$, and $s=k+ed \bmod n$: prove algebraically that the check $sG = R + eQ$ always holds for a correctly formed signature. Then explain why an attacker who does not know $d$ cannot produce a valid $(R,s)$ for a chosen $m$ without solving discrete log.
\end{exercise}

\begin{exercise}[Mining probability]
\label{ch36:ex-mining}
A proof-of-work target requires the double-SHA-256 of the header to start with $w$ zero \emph{bits}. Model each nonce as an independent trial. (a) What is the success probability $q$ per nonce? (b) What is the expected number of nonces tried before a success? (c) For $w=20$, roughly how many tries, and how does this compare to a real Bitcoin block needing on the order of $w=76$? (d) Using the geometric/exponential model, argue why the time since the last block tells you nothing about when the next one arrives.
\end{exercise}

\begin{exercise}[51\% scope]
\label{ch36:ex-fiftyone}
An attacker controls $60\%$ of a chain's hash power. For each of the following, state whether the attacker can do it and why: (a) reverse a transaction confirmed one block ago to double-spend their own coins; (b) spend coins sitting at an address whose private key they do not hold; (c) create $10{,}000$ new coins beyond the issuance schedule; (d) censor a specific user's transactions for as long as they hold the majority. Tie each answer to whether it is blocked by \emph{cryptography} or only by \emph{majority-following consensus}.
\end{exercise}

\begin{exercise}[What the math does not cover]
\label{ch36:ex-notcover}
For each loss scenario, name the failure and state whether the cryptography of this chapter was broken or merely bypassed: (a) a user types their twelve-word mnemonic into a phishing site; (b) malware swaps the destination address in the wallet UI just before the user signs; (c) a smart contract with a reentrancy bug is drained; (d) the coin's market price falls to zero. What single sentence summarizes the boundary these four share?
\end{exercise}

\begin{exercise}[Light-client trust]
\label{ch36:ex-spv}
An SPV phone wallet stores only $80$-byte headers and verifies payments via Merkle proofs. (a) Roughly how many bytes of headers for a chain of $800{,}000$ blocks? (b) What exactly does the wallet trust, and what does it verify itself? (c) Describe one thing an SPV client \emph{cannot} detect that a full node can (hint: think about whether a header's transactions obey the issuance and validity rules), and how committing to the full state (Ethereum's approach) narrows that gap.
\end{exercise}
